%% 段をまたぐフロート：最終ページで揃える，第3段で書いた図，3段にまたがる図，脚注
\documentclass[twocolumn,paper=b5,fontsize=9pt]{jlreq}
\usepackage[columns=3,debug,footnote=column]{evenend}
\setlength\columnseprule{0.4pt}
\newcommand\para{吾輩は猫である．名前はまだ無い．どこで生れたかとんと見当がつかぬ．何でも薄暗いじめじめした所でニャーニャー泣いていた事だけは記憶している．吾輩はここで始めて人間というものを見た．\par}
\newcommand\fig[2]{\centering\fbox{\parbox[c][#2][c]{.95\linewidth}{\centering #1}}\caption{#1}}
\begin{document}
\section{最終ページ}
\para\para\para\para\para\para\para\para\para\para\para\para\para\para\para\para\para\para\para\para\para\para\para\para\para\para\para\para
\begin{spanfigure}[b]{2}
\fig{第3段で書いた2段の図}{50pt}
\end{spanfigure}
\para\para\para
\begin{spanfigure}[t]{3}
\fig{3段にまたがる図}{40pt}
\end{spanfigure}
\para 脚注\footnote{脚注の本文．}を付ける．\para
\begin{spanfigure}[b]{2}
\fig{最終ページの2段の図}{45pt}
\end{spanfigure}
\para\para\para\para\para
\end{document}
